\centerline{\bf \Huge Пересчитаем ребра I}
\centerline{\bf \Huge Решения}

\problem{1}
В стране 1568 городов, из каждого выходит по 4 дороги. Сколько всего дорог в стране?

\textbf{Решение.} По лемме о рукопожатиях дорог будет $\dfrac{1568\cdot4}{2} = 3136$

\textit{Для ассистентов:} ...

\problem{2}
Можно ли изобразить на плоскости 9 отрезков так, чтобы каждый пересекал ровно 3 других? (считая, что в одной точке могут пересекаться не более двух отрезков)

\textbf{Решение.}
Рассмотрим граф, вершины которого – данные отрезки, а ребро соединяет две вершины тогда, когда два соответствующих отрезка пересекаются. У него 9 нечётных вершин, что противоречит лемме о рукопожатиях.

\textit{Для ассистентов:} помогите разобраться с тем, как тут найти граф. Что является вершиной, а что ребром.

\problem{3}
В некоторой стране 15 городов, каждый из которых соединен дорогами не менее чем с 7 другими. Докажите, что из любого города можно добраться до любого другого (возможно, проезжая через другие города).

\textbf{Решение.}
Рассмотрим два произвольных города и предположим, что они не соединены путем, то есть такой последовательностью дорог, в которой начало очередной дороги совпадает с концом предыдущей. Каждый из этих двух городов по условию соединен не менее, чем с семью другими; при этом все упомянутые города различны – ведь если какие-то два из них совпадают, то есть путь, соединяющий исходные города. Таким образом, мы насчитали не менее 16 городов. Противоречие.

\textit{Для ассистентов:} помогите дойти до того, что у любых двух вершин всегда есть одна общая.

\problem{4}
Каждый мальчик знаком с 7 мальчиками и 10 девочками, а каждая девочка — с 8 девочками и 9 мальчиками. Кого больше — мальчиков или девочек?

\textbf{Решение.}
Посчитаем количество дружб между девочками и мальчиками. Так как в каждой такой дружбе ровно один мальчик и он дружит с 10   девочками, то общее количество таких дружб равно 10$\cdot$количество мальчиков. С другой стороны, оно же равно 9$\cdot$количество девочек. Значит девочек больше, чем мальчиков.

\textit{Для ассистентов:} посчитайте вместе количество дружб

\problem{5}
У Элианы 5 друзей среди одноклассников. У остальных её одноклассников 4,6 или 8 друзей. И только у Эрдема всего один друг. Докажите, что Элиана может отправить Эрдему записку, если каждый будет передавать записку одному из своих друзей.

\textbf{Решение.} Пусть Элиана никак не сможет донести записку до Эрдема. Это значит, что Элиана и Эрдем находятся в разных компонентах связности с Эрдемом. Но в компоненте Элианы одна нечетная вершина и остальные чётные, что противоречит лемме о рукопожатиях.

\textit{Для ассистентов:} скорее всего ученику будет сложно оперировать понятиями "компонента связности". Поэтому придумайте какой-нибудь альтернативный способ, лучше всего нарисуйте какой-нибудь пример.

\problem{6}
Юноши и девушки пожимали друг другу руки (юноши — только девушкам, девушки — только юношам). Каждый юноша пожал 13 рук, а каждая девушка — 10 рук. При этом какие-то пары могли жать руки несколько раз. Докажите, что было не менее 130 рукопожатий.

\textbf{Решение.}
Пусть всего было $x$ юношей и $y$ девушек. Рассмотрим граф, в котором юношам и девушкам будут соответствовать вершины, а рукопожатиям — ребра.

Суммарная степень всех вершин-юношей равна $13 x$, а всех вершин-девушек — $10 y$. Меж тем, каждое ребро ведет от юноши к девушке, поэтому эти две суммы должны быть равны, более того, равны они общему числу ребер. Поэтому мы имеем равенство $13 x=10 y$. Отсюда $x:$ 10, значит, $x \geqslant 10$. Значит, общее количество ребер, равное $13 x$, не меньше $13 \cdot 10=130$, что и требовалось доказать.

